%% CV for Aalborg Universitet - Ph.d.-stipendiat (Kommunikations- og radarsystemer til UAS)
%% Compiled with: lualatex

\documentclass[11pt,a4paper,sans]{moderncv}
\moderncvstyle{banking}
\moderncvcolor{blue}

\renewcommand*{\firstnamestyle}[1]{{\fontsize{34}{36}\bfseries\upshape\color{color1}#1}}
\renewcommand*{\lastnamestyle}[1]{{\fontsize{34}{36}\bfseries\upshape\color{color1}#1}}
\renewcommand*{\sectionstyle}[1]{{\sectionfont\color{color1}#1}}

\usepackage[utf8]{inputenc}
\usepackage{hyperref}
\hypersetup{
    colorlinks=true,
    linkcolor=blue,
    filecolor=magenta,
    urlcolor=blue,
    pdftitle={Jacob El-Omar - CV},
    pdfpagemode=FullScreen,
}
\usepackage[scale=0.80]{geometry}
\usepackage{import}

\name{Jacob}{El-Omar}
\address{Aalborg, Denmark}{}{}
\phone[mobile]{+45 53 35 27 82}
\email{elomarjc@gmail.com}
\extrainfo{\href{https://linkedin.com/in/jacob-el-omar}{LinkedIn} \textbullet{} \href{https://github.com/elomarjc}{GitHub}}

\begin{document}

\makecvtitle

% ============================================================
%     PROFILE
% ============================================================

\vspace{6pt}
\small{Civilingeniør (MSc) i Elektroniske Systemer fra Aalborg Universitet med speciale i signalbehandling, stokastiske systemer, adaptiv filtrering og algoritmeudvikling i Python og MATLAB. Erfaring med stokastisk estimering, sensorintegration og hardware-software co-design samt tværfagligt projektsamarbejde. Motiveret for ph.d.-forskning inden for robuste mesh-kommunikations- og radarsystemer til droner i GNSS-benægtede miljøer ved Institut for Elektroniske Systemer (AAU).}

% ============================================================
%     TECHNICAL SKILLS
% ============================================================

\section{Faglige Kompetencer}
\vspace{1pt}
\begin{itemize}

\item \textbf{Trådløs Kommunikation \& Netværk}: Trådløse kanalmodeller (small- \& large-scale), PHY-evaluering, MIMO, Random/Multiple Access, Markov Chains, Discrete Event Simulation (TrueTime), CAN-busser samt Network Calculus.

\item \textbf{Signalbehandling \& Stokastik}: Adaptiv filtrering (LMS, NLMS, RLS), FastICA, stokastisk estimering, Kalman-filtrering, linkbudgetter og kanalmodellering.

\item \textbf{Algoritmeudvikling \& Simulering}: Python (NumPy, SciPy, Matplotlib), MATLAB / Simulink, C/C++, stokastisk modellering, dataset-generering og systememulering (Sionna).

\item \textbf{Embedded \& Reeltidssystemer}: Mikrocontroller-programmering (Cypress PSoC i C), I2C/IMU sensorintegration, reeltidsfiltrering og hardware-in-the-loop test.

\item \textbf{Systemdesign \& Robotik}: B-dot og LQR regulering, Ultra-Wideband (UWB) positionsbestemmelse, kollisionsforebyggelse samt Altium PCB \& Onshape konstruktion.

\item \textbf{Samarbejde \& Formidling}: Problembaseret læring (PBL), tværfagligt projektsamarbejde med eksterne forsknings- og industripartnere, teknisk dokumentation.

\end{itemize}

% ============================================================
%     EDUCATION
% ============================================================

\section{Uddannelse}
\vspace{1pt}

\cventry{Sep 2023--Jun 2025}{MSC i Elektroniske Systemer}{Aalborg Universitet}{Aalborg, Danmark}{}{
Kandidatspeciale: ``Adaptive Noise Cancellation For Electronic Stethoscopes.'' Reeltids-signalbehandling, adaptiv filtrering og stokastisk estimering i Python og MATLAB.
}

\vspace{3pt}

\cventry{Sep 2020--Aug 2023}{BSc i Elektronik og IT}{Aalborg Universitet}{Aalborg, Danmark}{}{
Specialisering i reguleringsteknik, digitale systemer og embedded C-programmering (Cypress PSoC, satellitsimulator CoM, MiR200 robotstyring).
}

% ============================================================
%     PROFESSIONAL EXPERIENCE
% ============================================================

\section{Erhvervserfaring}
\vspace{3pt}
\cventry{Sep 2024--Jun 2025}{AI/ML Engineer (Praktikant)}{Ai Health Highway}{Aalborg, Danmark / Remote}{}{
\begin{itemize}
    \item Udviklede og validerede signalbehandlings- og filtreringsalgoritmer i Python til støjreduktion i krævende miljøer.
    \item Implementerede adaptiv filtrering (LMS, NLMS, RLS) og FastICA til adskillelse af komplekse akustiske signaler.
    \item Samarbejdede med internationale forskningsteams om algoritmeoptimering og teknisk dokumentation.
\end{itemize}}

\vspace{3pt}

\cventry{Jan 2026--Nuværende}{IT- \& Webudvikler (Frivillig)}{Dawah Danmark}{Aalborg, Danmark}{}{
\begin{itemize}
    \item Vedligeholdelse af organisatorisk web- og databaseinfrastruktur (WordPress, MySQL) med fokus på oppetid og ydeevne.
\end{itemize}}

\vspace{3pt}

\cventry{Jun 2023--Aug 2024}{Lager- og Logistikmedarbejder}{Lyn Express}{Aalborg, Danmark}{}{
\begin{itemize}
    \item Håndtering af komplekse logistikarbejdsgange og dataopgørelser under højt arbejdstempo.
\end{itemize}}

% ============================================================
%     SELECTED PROJECTS
% ============================================================

\section{Udvalgte Projekter}
\vspace{1pt}
\begin{itemize}

\item{\textbf{AAUSAT6 ADCS Satellitregulering (AAU)}: Udvikling af B-dot detumbling-algoritmer og Kalman-filtrering i MATLAB/Simulink til positions- og orienteringsbestemmelse på CubeSat-platforme.}

\item{\textbf{UWB Robot Kollisionsforebyggelse (AAU)}: Design af prædiktivt navigations- og kollisionsforebyggelsessystem for MiR200 mobil robot baseret på Ultra-Wideband (UWB) sporing og regressionsmodeller.}

\item{\textbf{Center of Mass Kalibrering (AAU)}: Algoritmer til estimering af tyngdepunkt for 3-DoF satellitsimulator ved hjælp af kinematik og Kalman-filtrering i MATLAB.}

\end{itemize}

% ============================================================
%     LANGUAGES
% ============================================================

\section{Sprog}
\vspace{1pt}
\begin{itemize}
\item Dansk (Modersmål), Engelsk (Flydende i skrift og tale).
\end{itemize}

% ============================================================
%     REFERENCES
% ============================================================

\section{Referencer}
\vspace{1pt}
\begin{itemize}
\item Kan oplyses efter anmodning.
\end{itemize}

\end{document}
